\section{\textbf{Conclusão}}

A análise das estimativas de massa estelar no aglomerado das Plêiades, utilizando tanto o método de ajuste de isócronas quanto a relação massa-magnitude (KNN), permitiu uma comparação detalhada entre essas abordagens e os catálogos de referência disponíveis. Os resultados mostram que o método de relação massa-magnitude (KNN) foi mais consistente, especialmente para estrelas de massa intermediária, apresentando menor dispersão nas estimativas. Por outro lado, o ajuste de isócronas demonstrou limitações, particularmente para estrelas de maior massa, onde foram observadas subestimações. Essas discrepâncias ressaltam a importância de selecionar o método mais adequado de acordo com a faixa de massa e a população estelar em análise. As comparações realizadas com os dados dos catálogos de \citeonline{lodieu2019fiveD}, \citeonline{tess2019} e \citeonline{delfosse2000accurate} destacaram as diferenças na precisão dos métodos, reforçando que a escolha do catálogo e da metodologia impacta diretamente a confiabilidade das estimativas. Assim, para estudos futuros, o uso de modelos mais robustos, como a relação massa-magnitude, pode ser preferível para aglomerados como as Plêiades, onde as massas intermediárias são predominantes.